% Source: https://artofproblemsolving.com/community/c3h2483908_ant_and_a_rope
% Stable identity: aops-topic-2483908-q01
Describe the state of the rope at any moment by the set of pairs (position, direction) of the ants on it. When two ants meet at a point and both reverse, just before the meeting there is one ant at that point moving right and one moving left, and just after it the same is true. So the set of (position, direction) pairs evolves exactly as if the two ants had passed through each other without turning.

Consequently, imagine ``ghosts'': each ghost starts where an ant starts, with the same direction, and walks straight on at $1$ metre per second, ignoring all other ghosts, until it falls off an end. At every moment, the set of (position, direction) pairs of the ants equals that of the ghosts; only the names of the walkers differ. (Two ghosts moving in the same direction keep a constant distance apart, so at most two ghosts, one in each direction, are ever at the same point, and every meeting of ants is a meeting of exactly two of them.)

Each ghost walks straight towards an end that is at most $10$ metres away, so it falls off within $10$ seconds. Hence after $10$ seconds there are no ghosts on the rope, and therefore no ants on the rope either: all the ants have fallen off within $10$ seconds. $\blacksquare$

(The bound is sharp: a single ant starting at one end and walking towards the other end needs exactly $10$ seconds.)
